%% kapa: exam=1 round=22 year=2011 period=2 form=A subject=appraisal_law subject_ko=부동산관계법규 seq=1 number=41
%% source: Q-Net 감정평가사 기출문제 (cst00309) artlSeq=5036790 file=extracted/1008/5036790/2011년 제22회 감정평가사 1차시험 기출문제/2교시.hwp
%% conversion: kordoc 4.24.51 -> tools/gen1.py (auto); review=proofread-gpt-astra; flags=none
\begin{question}{41}
국토의 계획 및 이용에 관한 법령상 지방자치단체의 의견청취에 관한 설명으로 옳지 않은 것은?
\begin{choices}
  \item 시ㆍ도지사, 시장 또는 군수는 광역도시계획을 수립하거나 변경하려면 미리 관계 시ㆍ도, 시 또는 군의 의회와 관계 시장 또는 군수의 의견을 들어야 한다.
  \item 광역도시계획이 수립되는 시ㆍ도, 시 또는 군의 의회와 관계 시장 또는 군수는 특별한 사유가 없으면 30일 이내에 시ㆍ도지사, 시장 또는 군수에게 의견을 제시하여야 한다.
  \item 시장은 도시기본계획을 수립하거나 변경하려면 미리 그 시의회의 의견을 들어야 한다.
  \item 시장은 도시기본계획을 변경하려면 미리 공청회를 열어 주민과 관계 전문가 등으로부터 의견을 들어야 하며, 공청회에서 제시된 의견이 타당하다고 인정하면 도시기본계획에 반영하여야 한다.
  \item 시장 또는 군수가 녹지지역의 지정에 관한 도시관리계획을 입안할 경우에는 해당 지방의회의 의견청취를 생략할 수 있다.
\end{choices}
\end{question}
\answer{5}
